% Tipografía y página
\usepackage{fontspec}
\defaultfontfeatures{Ligatures=TeX}
\setmainfont{lmroman10-regular.otf}[
  BoldFont=lmroman10-bold.otf,
  ItalicFont=lmroman10-italic.otf,
  BoldItalicFont=lmroman10-bolditalic.otf
]
\setsansfont{lmsans10-regular.otf}[
  BoldFont=lmsans10-bold.otf,
  ItalicFont=lmsans10-oblique.otf,
  BoldItalicFont=lmsans10-boldoblique.otf
]
\setmonofont{lmmono10-regular.otf}[
  BoldFont=lmmonolt10-bold.otf,
  ItalicFont=lmmono10-italic.otf,
  Ligatures=NoCommon
]

\usepackage{babel}
\IfFileExists{spanish.ldf}{%
  \babelprovide[main,import=es]{spanish}%
}{%
  \babelprovide[main,import=es]{spanish}%
}
\setlocalecaption{spanish}{contents}{Índice general}
\setlocalecaption{spanish}{listtable}{Índice de tablas}
\setlocalecaption{spanish}{listfigure}{Índice de figuras}
\setlocalecaption{spanish}{table}{Tabla}
\setlocalecaption{spanish}{appendix}{Anexo}

\usepackage{amsmath}
\usepackage{unicode-math}
\setmathfont{latinmodern-math.otf}
\usepackage[left=27mm,right=26mm,top=25mm,bottom=27mm,headheight=15pt]{geometry}
\usepackage[protrusion=true,expansion=false]{microtype}
\usepackage{ragged2e}
\linespread{1.07}
\setlength{\parindent}{0pt}
\setlength{\parskip}{.48em plus .12em minus .08em}
\setlength{\emergencystretch}{1.3em}
\clubpenalty=7500
\widowpenalty=7500
\raggedbottom

% Paleta editorial
\usepackage[table]{xcolor}
\definecolor{Tinta}{HTML}{192C3D}
\definecolor{Azul}{HTML}{173B52}
\definecolor{Verde}{HTML}{086D65}
\definecolor{Gris}{HTML}{466074}
\definecolor{Linea}{HTML}{CAD9DF}
\definecolor{Papel}{HTML}{F2F7F8}
\definecolor{Suave}{HTML}{E4F2EF}

% Jerarquía y encabezados
\usepackage{titlesec}
\usepackage{fancyhdr}
\setcounter{secnumdepth}{1}
\setcounter{tocdepth}{1}
\newcommand{\EtiquetaCapitulo}{%
  \sffamily\bfseries\small\color{Verde}\MakeUppercase{\chaptertitlename}\enspace\thechapter}
\titleformat{\chapter}[display]
  {\sffamily\bfseries\Huge\color{Azul}\RaggedRight}
  {\EtiquetaCapitulo}{.7ex}{}[\vspace{.6ex}{\color{Linea}\titlerule[.65pt]}]
\titleformat{name=\chapter,numberless}[display]
  {\sffamily\bfseries\Huge\color{Azul}\RaggedRight}
  {}{0pt}{}[\vspace{.6ex}{\color{Linea}\titlerule[.65pt]}]
\titlespacing*{\chapter}{0pt}{1.3ex}{2.2ex}
\titleformat{\section}
  {\sffamily\bfseries\Large\color{Azul}\RaggedRight}
  {\textcolor{Verde}{\thesection}}{.7em}{}
\titlespacing*{\section}{0pt}{2.7ex plus .4ex}{.9ex}
\titleformat{\subsection}
  {\sffamily\bfseries\large\color{Azul}\RaggedRight}
  {}{0pt}{}
\titlespacing*{\subsection}{0pt}{1.9ex plus .4ex}{.55ex}

\newcommand{\NotaParte}{}
\titleformat{\part}[display]
  {\sffamily\bfseries\fontsize{38}{43}\selectfont\color{Azul}\RaggedRight}
  {\color{Verde}\normalfont\sffamily\bfseries\large\MakeUppercase{\partname}\enspace\thepart}
  {1.5ex}{}[%
    \par\vspace{.8ex}{\color{Verde}\rule{19mm}{1.1pt}}\par
    \vspace{1.6ex}%
    \normalfont\Large\color{Gris}\NotaParte]
\newcommand{\ParteLibro}[2]{\renewcommand{\NotaParte}{#2}\part{#1}}

\pagestyle{fancy}
\fancyhf{}
\fancyhead[L]{\sffamily\footnotesize\color{Azul}\TituloLibro}
\fancyhead[R]{\sffamily\footnotesize\color{Gris}\nouppercase{\leftmark}}
\fancyfoot[R]{\sffamily\footnotesize\color{Azul}\thepage}
\renewcommand{\headrule}{\vskip-1.5pt\hbox to\headwidth{\color{Linea}\leaders\hrule height.45pt\hfill}}
\renewcommand{\chaptermark}[1]{\markboth{#1}{}}
\fancypagestyle{plain}{%
  \fancyhf{}\fancyfoot[R]{\sffamily\footnotesize\color{Azul}\thepage}%
  \renewcommand{\headrulewidth}{0pt}%
  \renewcommand{\headrule}{}}

\usepackage{enumitem}
\setlist{topsep=.28em,itemsep=.2em,parsep=0pt,partopsep=0pt,leftmargin=1.65em}
\usepackage{needspace}
\usepackage[most]{tcolorbox}
\newtcolorbox{resultados}[1][Resultados de aprendizaje]{
  enhanced, breakable, colback=Suave, colframe=Suave, boxrule=0pt, arc=1.2mm,
  borderline west={1.1pt}{0pt}{Verde}, left=3.6mm, right=3.6mm, top=2.0mm,
  bottom=2.0mm, before skip=.9em, after skip=.9em,
  title={#1}, fonttitle=\sffamily\bfseries\small, coltitle=Verde,
  colbacktitle=Suave, toptitle=2mm, bottomtitle=0mm
}

% Cuadros, figuras, notas y referencias
\usepackage{graphicx}
\usepackage{caption}
\DeclareCaptionFont{rótulo}{\sffamily\bfseries\color{Azul}}
\captionsetup{font=small,labelfont=rótulo,labelsep=period,
  textfont=bf,justification=raggedright,singlelinecheck=false}
\captionsetup[table]{position=above,skip=5pt}
\captionsetup[figure]{position=below,skip=7pt}
\counterwithout{table}{chapter}
\counterwithout{figure}{chapter}
\usepackage{array}
\usepackage{booktabs}
\usepackage{longtable}
\newlength{\AnchoTabla}
\newcolumntype{L}[1]{>{\RaggedRight\arraybackslash}p{#1\AnchoTabla}}
\newcolumntype{R}[1]{>{\raggedleft\arraybackslash}p{#1\AnchoTabla}}
\newcolumntype{C}[1]{>{\centering\arraybackslash}p{#1\AnchoTabla}}
\newenvironment{tablalibro}[2]{%
  \par\needspace{5.3\baselineskip}\begingroup\small
  \setlength{\tabcolsep}{6pt}%
  \setlength{\AnchoTabla}{\dimexpr\linewidth-#1\tabcolsep-#1\tabcolsep+2\tabcolsep\relax}%
  \setlength{\LTleft}{0pt}\setlength{\LTright}{0pt}%
  \setlength{\LTcapwidth}{\linewidth}%
  \setlength{\LTpre}{.55em}\setlength{\LTpost}{.8em}%
  \renewcommand{\arraystretch}{1.24}%
  \arrayrulecolor{Linea}%
  \begin{longtable}{@{}#2@{}}%
}{\end{longtable}\endgroup}
\newcommand{\CabeceraTabla}[1]{\sffamily\bfseries\color{Azul}#1}
\newcommand{\NotaLibro}[1]{{\footnotesize\color{Gris}\textit{Nota.} #1}}
\newcommand{\SigueTabla}[1]{\multicolumn{#1}{@{}l@{}}{\footnotesize\color{Gris}Tabla \thetable\space (continuación)}\\}
\newcommand{\ContinuaTabla}[1]{\multicolumn{#1}{@{}r@{}}{\footnotesize\color{Gris}Continúa en la página siguiente.}\\}
\newcommand{\RefTabla}[1]{\hyperref[#1]{tabla~\ref*{#1}}}
\newcommand{\RefFigura}[1]{\hyperref[#1]{figura~\ref*{#1}}}
\newcommand{\RefCapitulo}[1]{\hyperref[#1]{capítulo~\ref*{#1}}}
\newcommand{\RefAnexo}[1]{\hyperref[#1]{anexo~\ref*{#1}}}

\usepackage{tikz}
\usetikzlibrary{arrows.meta,positioning,shapes.geometric,calc}
\tikzset{
  diagrama nodo/.style={font=\sffamily\small,align=center,inner sep=5pt,
    draw=Linea,fill=Papel,text=Azul,thick,minimum height=8mm},
  diagrama bloque/.style={diagrama nodo,rounded corners=1.3mm},
  diagrama proceso/.style={diagrama nodo},
  diagrama terminal/.style={diagrama nodo,ellipse,fill=Suave,draw=Verde,minimum width=16mm},
  diagrama datos/.style={diagrama nodo,trapezium,
    trapezium left angle=71,trapezium right angle=109,trapezium stretches=true,fill=Suave},
  diagrama decision/.style={diagrama nodo,diamond,aspect=2.45,fill=Suave,draw=Verde,
    inner sep=1.5pt},
  diagrama flecha/.style={-{Latex[length=2.2mm]},draw=Verde,line width=.76pt,rounded corners=1.3pt},
  diagrama union/.style={draw=Gris,line width=.55pt},
  diagrama rama/.style={font=\sffamily\footnotesize,text=Gris,fill=white,inner sep=2pt}
}
\NewDocumentCommand{\FiguraLibro}{m m m m}{%
  \par\addvspace{.8em}%
  \noindent\begin{minipage}{\linewidth}
    \centering
    \input{figuras/#1}\par
    \captionof{figure}{#2}\label{#3}%
    \if\relax\detokenize{#4}\relax\else
      \RaggedRight\NotaLibro{#4}\par
    \fi
  \end{minipage}\par\addvspace{.8em}%
}

% Listados de texto plano y código incluidos desde los archivos ejecutables.
\usepackage{listings}
\lstdefinestyle{libro}{
  basicstyle=\small\ttfamily\color{Tinta}, backgroundcolor=\color{Papel},
  frame=leftline, framerule=1pt, rulecolor=\color{Linea}, framesep=5pt,
  xleftmargin=2mm, xrightmargin=.75mm, aboveskip=.7em, belowskip=.65em,
  columns=fullflexible, keepspaces=true, upquote=true, tabsize=4,
  showstringspaces=false, breaklines=true, breakatwhitespace=false,
  keywordstyle=\color{Verde}\bfseries, commentstyle=\color{Gris}\itshape,
  stringstyle=\color{Azul},
  literate=
    {á}{{á}}1 {é}{{é}}1 {í}{{í}}1 {ó}{{ó}}1 {ú}{{ú}}1
    {Á}{{Á}}1 {É}{{É}}1 {Í}{{Í}}1 {Ó}{{Ó}}1 {Ú}{{Ú}}1
    {ñ}{{ñ}}1 {Ñ}{{Ñ}}1 {ü}{{ü}}1 {²}{{²}}1 {←}{{$\leftarrow$}}1
}
\NewDocumentCommand{\Python}{O{} m}{\lstinputlisting[style=libro,language=Python,#1]{recursos/#2}}
\NewDocumentCommand{\Pseudocodigo}{m}{\lstinputlisting[style=libro]{recursos/pseudocodigo/#1}}
\NewDocumentCommand{\Salida}{m}{\lstinputlisting[style=libro]{recursos/salidas/#1}}
\newcommand{\Codigo}[1]{\texttt{#1}}

\newcommand{\CapituloFinal}[1]{%
  \chapter*{#1}\markboth{#1}{}\addcontentsline{toc}{chapter}{#1}}
\newcommand{\SeccionLibro}[2][]{%
  \section*{#2}%
  \if\relax\detokenize{#1}\relax
    \addcontentsline{toc}{section}{#2}%
  \else
    \addcontentsline{toc}{section}{#1}%
  \fi
}
\newcounter{ejemplo}[chapter]
\renewcommand{\theejemplo}{\thechapter.\arabic{ejemplo}}
\newcommand{\Ejemplo}[1]{%
  \refstepcounter{ejemplo}\subsection*{Ejemplo \theejemplo. #1}}

\usepackage[authoryear,round]{natbib}
\setlength{\bibhang}{1.5em}
\setlength{\bibsep}{.45em}
\renewcommand{\bibsection}{\CapituloFinal{Referencias bibliográficas}}
\usepackage[unicode,bookmarksnumbered,bookmarksopen,
  colorlinks,linkcolor=Azul,citecolor=Verde,urlcolor=Verde]{hyperref}
\usepackage{bookmark}
\hypersetup{
  pdftitle={\TituloLibro. \SubtituloLibro},
  pdfsubject={\TipoManuscrito},
  pdflang={es-PE}
}
\AtBeginDocument{\color{Tinta}}
